\documentclass[a4paper,12pt]{article}
\usepackage[utf8]{inputenc}
\usepackage[T1]{fontenc}
\usepackage[english]{babel}
\usepackage{geometry}
\geometry{left=2.5cm,right=2.5cm,top=2.5cm,bottom=2.5cm}

\usepackage{lmodern}
\usepackage{xcolor}
\definecolor{titleblue}{RGB}{0,51,140}

\usepackage{hyperref}

\pagestyle{plain}

\begin{document}

\begin{center}
    {\Large\bfseries Statement of Research Interests}\\[6pt]
    \textbf{Candidate: Ismail AISSAOUI} \\
    PhD Position: Structural Generalization in Transformer-based LLMs \\
    Saarland University, Germany
\end{center}

\bigskip

\section*{Research Motivation}
My research interests lie at the intersection of structural inductive biases and the generalization capabilities of deep sequence models. During my engineering degree at ESI-SBA, I have been fascinated by how different architectural constraints (Transformers vs. State Space Models vs. GNNs) influence a model's ability to generalize to unseen structural patterns—a challenge known as "Structural Generalization."

\section*{Current Research \& Findings}
In my current project at Sonatrach, I am investigating the use of \textbf{Mamba-SSM (State Space Models)} and \textbf{xLSTM} as generative surrogates for industrial digital twins. My key finding is that while Transformers excel at capturing global context, they often struggle with the "Structural Latency Paradox" in real-time CPS environments. By implementing a dual-track architecture (Cloud/Edge), I explored how physics-informed constraints (\textbf{PINNs}) can act as a structural bias to ensure thermodynamic consistency in out-of-distribution scenarios.

Furthermore, my work on \textbf{Graph Neural Networks (GNNs)} for the Career Connect project allowed me to explore structural generalization in a non-Euclidean space. I researched how network centrality measures can be used to improve node representation stability, which I believe is directly applicable to the "Structural Generalization" problem in LLMs, where the model must handle recursive or long-range dependencies that follow specific structural rules.

\section*{Future Goals at Saarland University}
The PhD position at Saarland University offers a unique opportunity to formalize these empirical observations into a rigorous theoretical framework. I am particularly interested in:
\begin{enumerate}
    \item \textbf{Structural Inductive Biases}: Investigating how attention mechanisms can be modified to explicitly respect structural priors (e.g., hierarchy, recursion).
    \item \textbf{Generalization beyond Sequence Length}: Studying how Transformer-based LLMs can be made to generalize to structural complexities that exceed their training window.
    \item \textbf{Hybrid Architectures}: Exploring the synergy between Transformer-based LLMs and structured representations (like GNNs or SSMs) to enhance robustness.
\end{enumerate}

\section*{Conclusion}
I am eager to bring my background in AI engineering and my passion for structural modeling to the Saarland University research team. I am confident that my combination of algorithmic understanding and deep learning expertise will allow me to contribute meaningfully to the project.

\bigskip

\begin{flushright}
    \textbf{Ismail AISSAOUI}
\end{flushright}

\end{document}
